\documentclass[11pt]{article}
\usepackage[margin=1in]{geometry}
\usepackage{booktabs}
\usepackage{array}
\usepackage{hyperref}
\usepackage{float}

\title{RAG Assignment: Baseline vs.\ Cross-Encoder Re-Ranking}
\author{Replace with your name and roll number}
\date{19 October 2026}

\begin{document}
\maketitle

\section*{Setup}
Document: Gloomhaven rulebook (74 pages, scanned; text obtained with local OCR), split into 183 chunks of about 200 words with 40 words of overlap. Bi-encoder: \texttt{all-MiniLM-L6-v2} with cosine similarity over a NumPy matrix. Cross-encoder: \texttt{cross-encoder/ms-marco-MiniLM-L-6-v2}. LLM: \texttt{gemini-3.5-flash-lite}, temperature 0, identical prompt for both pipelines. Code: \url{https://github.com/murthymetta2002/RAG-assignment}.

\section{Architecture Comparison}
A bi-encoder encodes the query and each chunk separately into one vector each, and relevance is the cosine similarity between the two vectors. Chunk vectors can be computed once and stored, so search is fast, but each vector has to summarise the whole chunk without knowing the query. A cross-encoder takes the query and the chunk together as a single input, so self-attention can compare every query token with every chunk token, and it outputs a relevance score directly. This catches things like negation and ``except when'' wording that a single averaged vector blurs, but nothing can be precomputed.

\section{Qualitative and Rank-Shift Analysis}
Gold chunks were identified by a phrase from the rulebook that contains the true answer. Ranks refer to the baseline top 10 and the re-ranked top 10.

\begin{table}[H]
\centering
\small
\begin{tabular}{lp{6cm}lcc}
\toprule
Category & Query & Gold chunk & Rank before & Rank after \\
\midrule
keyword trap & Which status makes a figure lose one life at the start of its turn? & C062, C063 & 6 & 1 \\
buried & Which positive effects can a poisoned monster still receive? & C063 & 6 & 2 \\
exception & Must the party always go to Gloomhaven after a scenario? & C133 & 5 & 2 \\
exception & What happens to retaliate if the retaliating figure is killed? & C069 & 5 & 1 \\
negation & Can loot left on the map be collected after the scenario ends? & C117 & 5 & 1 \\
negation & Can the party lose more reputation than they have? & C119 & 7 & 4 \\
control & Can a figure stand inside a wall line? & C029 & 1 & 4 \\
control & Is a road event drawn when the scenario is in casual mode? & C021 & 2 & 4 \\
\bottomrule
\end{tabular}
\caption{Rank of the gold chunk before and after re-ranking.}
\end{table}

\noindent\textbf{Where re-ranking helped.} In five of the eight queries the gold chunk moved from rank 5--6 to rank 1--2, which is inside the top 3 given to the LLM.
\begin{itemize}
\item \emph{Keyword trap} (wound): the query says ``lose one life'' but the rulebook says ``suffers 1 damage''. The gold chunk moved from \#6 to \#1. Answer 1 said the context does not mention such a status; Answer 2 correctly named Wound.
\item \emph{Exception} (City Phase): the chunk with the ``except in the following situations'' clause moved from \#5 to \#2. Answer 1 gave a partly correct answer but cited the wrong rule for losing a scenario; Answer 2 listed the replay and force-linked exceptions correctly.
\item \emph{Exception} (retaliate): \#5 to \#1. Answer 1 said it could not answer; Answer 2 gave the rule that the retaliate bonus does not occur if the figure is killed.
\item \emph{Negation} (loot): \#5 to \#1. Answer 1 said the context had no answer; Answer 2 correctly said tokens and treasure left on the map cannot be looted.
\item \emph{Buried} (poison): \#6 to \#2, but the question was poorly phrased and Answer 2 still did not give a direct answer. The retrieval improved without the answer improving much.
\end{itemize}
In none of these cases did the baseline invent a rule. It omitted the key clause or refused to answer because the critical chunk was outside the top 3.

\noindent\textbf{Where re-ranking did not help.}
\begin{itemize}
\item Reputation query: \#7 to \#4, so the gold chunk still missed the top 3. Both answers were similar.
\item Wall line query: the gold chunk went from \#1 to \#4. Answer 1 was correct, Answer 2 said the context had no answer. The cross-encoder preferred chunks about walls in general over the one with the specific rule.
\item Casual mode query: \#2 to \#4. Both answers were still correct because other chunks in the top 3 repeated the rule.
\end{itemize}
So the cross-encoder improves results on average here but is not strictly better; it was trained on web-search passages (MS MARCO), not rulebooks.

\section{Latency Profiling}
Mean over 8 queries with 5 repeats each, after a warm-up run, on a CPU-only laptop. LLM time is excluded because it is the same for both pipelines.

\begin{table}[H]
\centering
\begin{tabular}{lrr}
\toprule
Stage & Mean (ms) & Std (ms) \\
\midrule
Baseline retrieval (embed query + top-10) & 15.7 & 2.2 \\
Cross-encoder re-rank (10 pairs) & 557.9 & 195.7 \\
Total for pipeline 2 & 573.6 & -- \\
\bottomrule
\end{tabular}
\caption{Retrieval latency per query.}
\end{table}

Re-ranking adds roughly 35 times the baseline retrieval time for these 10 candidates.

\section{Scaling Question}
A cross-encoder needs a full forward pass for every (query, chunk) pair and nothing can be cached, because the score depends on the query. Running it on a whole collection of $N$ chunks would cost $N$ forward passes per query, whereas a bi-encoder needs one forward pass for the query and a fast vector search over precomputed embeddings. Here $N=183$ and 10 pairs already cost about 560\,ms, so for a collection with millions of chunks the cross-encoder alone would be far too slow. That is why we use the bi-encoder to cut the collection down to a small candidate set first and only re-rank those.

\end{document}
